\section*{الفصل \ref{ch:b3:exoplanets} --- الكواكب الخارجية}

\begin{solution}{exo:b3:exoplanets:1}
(أ)~$K = (2\pi G/P)^{1/3}m_p/\bigl((M_\odot + m_p)^{2/3}\sqrt{1 - e^2}\bigr)$.
المشتري: $P = \qty{3.742e8}{s}$، $(2\pi G/P)^{1/3} = \num{1.039e-6}$
بوحدات SI، $K = \num{1.039e-6}\times\num{1.898e27}/(\num{1.990e30})^{2/3}/0.9988
= \qty{12.5}{m.s^{-1}}$. أمّا الأرض: $P = \qty{3.155e7}{s}$، $K =
\qty{0.089}{m.s^{-1}}$، نحو \qty{9}{cm.s^{-1}}. (ب)~\qty{3.5}{km/h} $=
\qty{0.97}{m.s^{-1}}$ و\qty{0.35}{km/h} $= \qty{0.097}{m.s^{-1}}$.
تبلغ إشارة المشتري ثلاثة عشر ضعف الدقّة الأولى؛ أمّا إشارة الأرض فتقع
دون الدقّة الثانية، وتضاف فوقها ضوضاء النجم نفسه.
\end{solution}

\begin{solution}{exo:b3:exoplanets:2}
(أ)~$(71492/695700)^2 = \num{0.0106}$، بانحدار مقداره \qty{1.06}{\percent}؛
$(6378/695700)^2 = \num{8.4e-5}$، أي 84 جزءًا في المليون. (ب)~$R_\star/a$:
الأرض $695700/\num{1.496e8} = \num{4.7e-3}$، أي نحو فرصة واحدة من 215؛
المشتري $695700/\num{7.785e8} = \num{8.9e-4}$، فرصة واحدة من 1100.
\end{solution}

\begin{solution}{exo:b3:exoplanets:3}
$m_p\sin i \approx K M_\odot^{2/3}(P/2\pi G)^{1/3}$ مع $P =
\qty{3.655e5}{s}$: $(P/2\pi G)^{1/3} = \num{9.55e4}$ بوحدات SI، ومنه $m_p
\sin i = 59\times(\num{1.9884e30})^{2/3}\times\num{9.55e4} =
\qty{8.9e26}{kg}$، أي 0.47 من كتلة المشتري أو 149 كتلة أرضية
(إنّ تصحيح $m_p$ في $M_\star + m_p$ أدنى من جزء واحد في الألف).
$a = \bigl(GM_\odot P^2/4\pi^2\bigr)^{1/3} = \qty{7.66e9}{m} =
\qty{0.051}{au}$.
\end{solution}

\begin{solution}{exo:b3:exoplanets:4}
(أ)~$4GM/c^2 = 4\times\num{6.674e-11}\times\num{9.94e29}/(\num{2.998e8})^2 =
\qty{2953}{m}$; $(D_S - D_L)/(D_LD_S) = 1/D_S\times(D_S - D_L)/D_L =
1/(\qty{8}{kpc}) = 1/\qty{2.469e20}{m}$. $\theta_\mathrm E =
\sqrt{2953/\num{2.469e20}} = \qty{3.46e-9}{rad} = \qty{0.71}{mas}$. عند
العدسة، $R_\mathrm E = \theta_\mathrm E D_L = \num{3.46e-9}\times
\qty{1.234e20}{m} = \qty{4.27e11}{m} = \qty{2.85}{au}$. (ب)~$A(1) =
3/\sqrt5 = 1.34$; $A(0.5) = 2.25/(0.5\times2.062) = 2.18$; $A(0.1) =
2.01/(0.1\times2.002) = 10.0$. (ج)~$t_\mathrm E = \num{4.27e11}/\num{2e5}
= \qty{2.13e6}{s}$، أي نحو 25 يومًا: أحداث \omterm{def:b3:exoplanets:microlensing}{التعديس الجذبوي الصغري} تدوم أسابيع.
\end{solution}

\begin{solution}{exo:b3:exoplanets:5}
(أ)~$a = (GM_\odot P^2/4\pi^2)^{1/3} = \qty{6.09e9}{m}$ (\qty{0.0407}{au})،
ومنه $a/R_\star = 8.75$. (ب)~مع $P/\pi = \qty{22.9}{h}$، ومن أجل $b = 0$:
\[
T_{14} = \frac P\pi\arcsin\frac{1.12}{8.75} = \qty{2.94}{h}, \qquad
T_{23} = \frac P\pi\arcsin\frac{0.88}{8.75} = \qty{2.31}{h}.
\]
من أجل $b = 0.7$: $\cos i = 0.7/8.75$، $i = \ang{85.4}$، $(a/R_\star)\sin i =
8.726$،
و\[
T_{14} = \frac P\pi\arcsin\frac{\sqrt{1.2544 - 0.49}}{8.726} = \qty{2.30}{h},
\qquad
T_{23} = \frac P\pi\arcsin\frac{\sqrt{0.7744 - 0.49}}{8.726} = \qty{1.40}{h}.
\]
(ج)~يستغرق الدخول والخروج $(T_{14} - T_{23})/T_{14} = \qty{21}{\percent}$
من العبور حين يكون مركزيًا، ويبلغ \qty{39}{\percent} عند $b = 0.7$؛ أمّا
العبور المماسّي، $b \to 1$، فلا يملك قاعًا مستويًا البتّة، ويُشبه حرف V.
\end{solution}

\begin{solution}{exo:b3:exoplanets:6}
$\bar\rho = M/\tfrac43\pi R^3$، مع $\tfrac43\pi R_\oplus^3 =
\qty{1.087e21}{m^3}$: ترابيست-1~b $1.374\times\num{5.97e24}/(1.116^3\times
\num{1.087e21}) = \qty{5.43}{g.cm^{-3}}$؛ كبلر-10~b
\qty{5.60}{g.cm^{-3}}؛ واسب-39~b \qty{0.167}{g.cm^{-3}}؛ \mbox{\textdir TLT HD~209458~b}
\qty{0.337}{g.cm^{-3}}. يماثل الكوكبان الأوّلان الأرض في الكثافة
(\qty{5.49}{g.cm^{-3}} بالثوابت نفسها) وصخريان؛ أمّا الكوكبان الأخيران
فهما عملاقان غازيان أخفّ من الماء، واسب-39~b أخفّ بثماني مرات من
المشتري.
\end{solution}

\begin{solution}{exo:b3:exoplanets:7}
القيمتان القصويان هما $\gamma + K(1 + e\cos\omega)$ و$\gamma - K(1 -
e\cos\omega)$. فرقهما $2K = \qty{60.0}{m.s^{-1}}$، ومنه $K =
\qty{30.0}{m.s^{-1}}$؛ ومتوسطهما $\gamma + Ke\cos\omega =
\qty{20.4}{m.s^{-1}}$، إذن $Ke\cos\omega = \qty{8.0}{m.s^{-1}}$
و$e\cos\omega = 0.27$. في المدار اللامركزي يمكث النجم قرب
الأوج، ويندفع عبر الحضيض، فيكون المنحنى غير متماثل،
ومنتصف قيمتيه القصويين لا يساوي متوسطه الزمني $\gamma$.
\end{solution}

\begin{solution}{exo:b3:exoplanets:8}
(أ)~$g = GM/R^2 = \num{6.674e-11}\times\num{5.332e26}/(\num{9.144e7})^2 =
\qty{4.26}{m.s^{-2}}$; $H = kT/\bar m g = \num{1.381e-23}\times1100/(2.3
\times\num{1.661e-27}\times4.26) = \qty{934}{km}$. ومع $R_\star =
\qty{6.533e8}{m}$، ويكون العمق $(\num{9.144e7}/\num{6.533e8})^2 =
\num{0.0196}$. (ب)~$2HR_p/R_\star^2 = 2\times\num{9.34e5}\times
\num{9.144e7}/(\num{6.533e8})^2 = \num{4.0e-4}$، أي 400 جزءًا في المليون؛
ارتفاعان سلّميان، أي 800: العملاق الغازي المنخفض الجاذبية هو الهدف الأسهل
\omterm{def:b3:exoplanets:transmission}{لمطيافية النفاذ}.
\end{solution}

\begin{solution}{exo:b3:exoplanets:9}
$p = 1.116\times6378/(0.1192\times695700) = 0.0858$, $p^2 = \num{7.37e-3}$.
$B_\lambda \propto 1/(e^{hc/\lambda kT} - 1)$ مع $hc/\lambda k =
\qty{959}{K}$ عند \qty{15}{\micro m}: النجم $1/(e^{0.374} - 1) = 2.20$؛
جانب النهار عند \qty{400}{K}، $1/(e^{2.40} - 1) = 0.0998$، النسبة \num{0.0453}،
العمق \num{3.3e-4} (334 جزءًا في المليون)؛ عند \qty{500}{K}، $1/(e^{1.92} -
1) = 0.172$، النسبة \num{0.078}، العمق \num{5.7e-4}. مئة كلفن
تغيّر العمق بمقدار \qty{70}{\percent}: العمق المقيس يحدّد درجة حرارة
جانب النهار بدقّة، وجانب النهار الحار يعني أنّ
قدرًا قليلًا من الحرارة ينتقل إلى جانب الليل.
\end{solution}

\begin{solution}{exo:b3:exoplanets:10}
(أ)~$S \propto L/d^2$، ومنه $d = \sqrt{0.005}\,\unit{au} = \qty{0.0707}{au}$؛
$P = 2\pi\sqrt{a^3/GM} = \qty{1.33e6}{s} = \qty{15.4}{d}$.
(ب)~$K = (2\pi G/P)^{1/3}M_\oplus/(0.2M_\odot)^{2/3} = \qty{0.75}{m.s^{-1}}$،
أي 8.4 ضعف تأثير الأرض في الشمس: يتحرّك النجم الخفيف أكثر،
وتساعد الدورة القصيرة أيضًا. (ج)~$(6378/139140)^2 = \num{2.1e-3}$، أي 25 ضعفًا
لعبور الأرض أمام الشمس. كلتا الإشارتين أكبر، وتقع العبورات
كلّ أسبوعين، واحتمال العبور $R_\star/a =
\num{0.013}$ يبلغ ثلاثة أضعاف نظيره للأرض: تمنح النجوم الصغيرة الكواكب الصخرية في
نطاقاتها الصالحة للحياة فرصةً لاكتشافها ودراستها.
\end{solution}

\begin{solution}{exo:b3:exoplanets:11}
باعتبار $r\,dr = -\mu\,d\mu$، تكون الشدّة المتوسطة $\frac1\pi\int_0^1
I\,2\pi r\,dr = 2\int_0^1 I(\mu)\mu\,d\mu = 2\bigl[\tfrac12 - u_1(\tfrac12 -
\tfrac13) - u_2\int_0^1(1 - \mu)^2\mu\,d\mu\bigr]$. التكامل الأخير هو
$\tfrac12 - \tfrac23 + \tfrac14 = \tfrac1{12}$، ومنه يكون المتوسط $1 - u_1/3 -
u_2/6$، وهو هنا $0.825$. يحجب كوكب صغير $\pi p^2 I(\mu)$ من أصل $\pi\times
0.825$: العمق $p^2 I(\mu)/0.825$. عند $r = 0$، $I = 1$: $1.212\,p^2$؛ عند
$r = 0.5$ ($\mu = 0.866$)، $I = 0.942$: $1.142\,p^2$؛ وعند $r = 0.9$ ($\mu =
0.436$)، $I = 0.695$: $0.842\,p^2$. سيقرأ نموذج القرص المنتظم
$1.212\,p^2$ على أنّه $p^2$، ويبالغ في تقدير نصف القطر بمقدار $\sqrt{1.212} = 1.10$،
\qty{10}{\percent}.
\end{solution}

\begin{solution}{exo:b3:exoplanets:12}
(أ)~$R_\odot/a$: \num{9.3e-3} عند \qty{0.5}{au}، \num{6.6e-3} عند
\qty{0.7}{au}، \num{4.7e-3} عند \qty{1}{au}. (ب)~$42000\times0.11\times
\num{6.6e-3} = 31$ كوكبًا يعبر. (ج)~لم يُعثر إلا على نحو ثلثها
ذلك أنّ بقيّتها أعطت عبورات ضحلة جدًّا (84 جزءًا في المليون
لكوكب كالأرض) أو قليلة جدًّا (واحد أو اثنان في السنة) فلم ترتفع فوق الضجيج. وهذه
النسبة هي كفاءة الكشف؛ ولو لم نصحّح بها، ولاحتمال العبور،
لكانت الأعداد تقلّل من تقدير مدى شيوع مثل هذه
الكواكب بعامل يقارب 460.
\end{solution}

\begin{solution}{pb:b3:exoplanets:1}
\textbf{1.} $M_\star = 0.80\times\num{1.9884e30} = \qty{1.591e30}{kg}$, $P =
\qty{1.0368e6}{s}$: $a = (GM_\star P^2/4\pi^2)^{1/3} = \qty{1.4245e10}{m}
= \qty{0.0952}{au}$.

\textbf{2.} $v = 2\pi a/P = \qty{86.3}{km.s^{-1}}$.

\textbf{3.} $m_p\sin i = KM_\star^{2/3}(P/2\pi G)^{1/3} = 2.59\times\num{1.363e20}
\times\num{1.352e5} = \qty{4.77e25}{kg}$، \num{7.99} كتلة أرضية.

\textbf{4.} $a_\star = a\,m_p/M_\star = \num{1.4245e10}\times\num{4.77e25}/
\num{1.591e30} = \qty{427}{km}$; $2\pi a_\star/P = \qty{2.59}{m.s^{-1}}$,
وهي $K$ المرصودة.

\textbf{5.} $2.59$ ضعفًا لدقّة قدرها \qty{1}{m.s^{-1}}، و26 ضعفًا لواحدة منها قدرها
\qty{0.1}{m.s^{-1}}: الكوكب في متناول أفضل الأجهزة،
وبالكاد في متناول الجيل السابق.

\textbf{6.} $m_p\sin i \propto K\sqrt{1 - e^2}$: $\times0.980$، 7.83 كتلة أرضية
\qty{2}{\percent} أقلّ، أي أقلّ بنسبة \qty{2}{\percent}.

\textbf{7.} $p = \sqrt{\num{8.38e-4}} = 0.02895$.

\textbf{8.} $R_\star = 0.76\times\qty{695700}{km} = \qty{528700}{km}$; $R_p =
\qty{15310}{km} = 2.40$ نصف قطر أرضي.

\textbf{9.} $a/R_\star = \num{1.4245e10}/\num{5.287e8} = 26.94$.

\textbf{10.} $T_{14} = (\qty{288}{h}/\pi)\arcsin(1.02895/26.94) = \qty{3.50}{h}$.

\textbf{11.} $\sqrt{(1 + p)^2 - b^2} = (a/R_\star)\sin(\pi T_{14}/P) = 26.94\times
\sin(\pi\times3.23/288) = 0.949$ (مع $\sin i \approx 1$)؛ $b^2 =
1.0587 - 0.9008 = 0.158$, $b = 0.40$.

\textbf{12.} $\cos i = b\,R_\star/a = 0.40/26.94 = 0.0148$, $i = \ang{89.15}$.

\textbf{13.} $R_\star/a = 1/26.94 = 0.037$: فرصة واحدة من 27.

\textbf{14.} $\sin i = 0.99989$: الكتلة الحقيقية هي 7.99 كتلة أرضية، \omterm{def:b3:exoplanets:rv}{والكتلة الدنيا}
بفارق لا يتجاوز \num{1e-4}.

\textbf{15.} $\bar\rho = \num{4.77e25}/\bigl(\tfrac43\pi(\num{1.531e7})^3\bigr) =
\num{4.77e25}/\num{1.502e22} = \qty{3180}{kg.m^{-3}}$,
\qty{3.18}{g.cm^{-3}}.

\textbf{16.} الكوكب أكبر بمقدار $2.40/1.74 = 1.38$ مرّة من كوكب صخري له كتلته،
وحجمه أكبر بمقدار 2.6 مرّة. لا يكفي الصخر وحده لتكوينه: إنّه يحتاج إلى
غلاف من غاز خفيف، بضعة في المئة من كتلته من الهيدروجين والهليوم،
أو إلى نسبة كبيرة من الماء.

\textbf{17.} عند 2.40 نصف قطر أرضي يقع الكوكب فوق الوادي، في قمّة 2--3 نصف قطر أرضي
القمّة: \omterm{def:b3:exoplanets:classes}{ما دون نبتون}.

\textbf{18.} $g = Gm_p/R_p^2 = \num{6.674e-11}\times\num{4.77e25}/(\num{1.531e7})^2 =
\qty{13.6}{m.s^{-2}}$.

\textbf{19.} $L = 4\pi R_\star^2\sigma T^4 = 4\pi(\num{5.287e8})^2\times\num{5.670e-8}
\times5000^4 = \qty{1.245e26}{W}$، 0.325 إشعاعية شمسية.

\textbf{20.} $S/S_\oplus = 0.325/0.0952^2 = 35.9$: بعيدًا جدًّا داخل الحافة الداخلية
للنطاق الصالح للحياة.

\textbf{21.} $L/4\pi a^2 = \qty{4.88e4}{W.m^{-2}}$; $T_\mathrm{eq} =
\bigl(\num{4.88e4}\times0.7/(4\sigma)\bigr)^{1/4} = \qty{623}{K}$.

\textbf{22.} $H = kT/\bar m g = \num{1.381e-23}\times623/(2.3\times\num{1.661e-27}
\times13.6) = \qty{166}{km}$.

\textbf{23.} $2HR_p/R_\star^2 = 2\times\num{1.66e5}\times\num{1.531e7}/
(\num{5.287e8})^2 = \num{1.82e-5}$، 18 جزءًا في المليون؛ وحزمةٌ بارتفاع
مقياسين للارتفاع، أي 36 جزءًا في المليون.

\textbf{24.} $H$ أصغر بمقدار $18/2.3 = 7.8$ مرّة، أي \qty{21}{km}: 2.3 جزءًا في
المليون لكلّ مقياس ارتفاع، و4.6 للحزمة. عند عشرة أجزاء في المليون، يُظهر الغلاف الجوي الغني بالهيدروجين
حزمة الماء عند أكثر من ثلاثة أضعاف
الضجيج؛ ويبدو الغلاف الجوي البخاري مسطّحًا، شأنه شأن غلاف غائم أو كوكب
بلا هواء.

\textbf{25.} \textbf{\qty[text-series-to-math=true]{3.2}{g/cm^3}}.
\end{solution}
